\documentclass[11pt,a4paper]{tubaf-report}
\usepackage[sans]{tubaf-fonts}
\SetKeys[tubaf-title]{black=false}

\usepackage[english]{babel}
\usepackage[utf8]{inputenc}
\usepackage[T1]{fontenc}
\usepackage{textcomp}
\usepackage{graphicx}
\usepackage{xurl}
\usepackage[
    colorlinks,
    citecolor=black,
    urlcolor=blue,
    linkcolor=black,
    bookmarksopen=true,
    hypertexnames=true
]{hyperref}
\usepackage{xcolor}
\usepackage{listings}

\definecolor{codebackground}{RGB}{245,247,250}
\definecolor{codeframe}{RGB}{190,195,200}
\definecolor{codecomment}{RGB}{0,110,70}

\lstdefinestyle{terminal}{
    backgroundcolor=\color{codebackground},
    basicstyle=\ttfamily\small,
    frame=single,
    rulecolor=\color{codeframe},
    breaklines=true,
    breakatwhitespace=false,
    columns=fullflexible,
    keepspaces=true,
    showstringspaces=false,
    upquote=true,
    xleftmargin=0.5em,
    framexleftmargin=0.5em
}

\begin{document}

\title[OpenGeoSys Installation Manual]{OpenGeoSys Training Course}
\subtitle{Numerical Modelling of Coupled Thermo-Hydro-Mechanical Processes in Geological Media}
\author[Mehran Ghasabeh]{Mehran Ghasabeh}
\date{October 2026}
\subject{Installation and Course Environment Manual}
\publishers{Institute of Geotechnics, TU Bergakademie Freiberg}
\moretitle{Course contact: \href{mailto:mehran.ghasabeh@ifgt.tu-freiberg.de}{mehran.ghasabeh@ifgt.tu-freiberg.de}}
\titlefoot{Joint OpenGeoSys Training Course -- UFZ, KIGAM, Zhejiang University, and TUBAF}
\setlength{\logoheadsize}{2.5cm}
\makeatletter
\renewcommand{\tubaf@title@logo}{\insertlogo{5.5cm}}
\makeatother
\logo{\includegraphics{UFZ_KIGAM_ZJU}}

\maketitle

\tableofcontents
\clearpage

\chapter{Course Overview}

\noindent
In geosciences and related engineering fields, computational geomechanics plays a critical role in understanding complex geological processes and addressing a wide range of geotechnical and geoscientific applications. These processes often involve multiphysical coupling across different spatial and temporal scales, requiring robust numerical frameworks capable of addressing inherent nonlinearities and complex material behaviour. Accurate numerical treatment of these coupled processes is therefore essential for stable, efficient, and reliable simulations of geomechanical systems.


The course introduces the fundamental concepts and practical application of numerical methods for solving geomechanical problems using OpenGeoSys, an open-source finite element software platform for the simulation of coupled thermo-hydro-mechanical-chemical (THMC) processes in geosciences and related engineering fields.


\noindent
The course will cover:
\begin{itemize}
    \item Introduction to numerical modelling in geomechanics,
    \item Fundamentals of the finite element method for geomechanical problems,
    \item Introduction to the OpenGeoSys framework and simulation workflow,
    \item Model setup: geometry, mesh, material properties, and initial/boundary conditions,
    \item Constitutive models for geomechanical simulations,
    \item Numerical simulation of deformation and stress in geomaterials,
    \item Coupled hydro-mechanical and thermo-hydro-mechanical processes,
    \item Practical examples and hands-on simulations using OpenGeoSys,
    \item Post-processing and interpretation of numerical results.
\end{itemize}

The course combines theoretical explanations with practical exercises, enabling participants to develop and run numerical models using OpenGeoSys and to gain practical experience in the simulation of geomechanical and coupled THMC processes.


The course is intended for researchers, engineers, graduate students, and professionals working in geotechnical engineering, geomechanics, geological engineering, computational geosciences, and related fields.


\chapter{OpenGeoSys Installation and Requirements on Ubuntu}
\label{chap:ubuntu-installation}

These instructions use Python~3.13 on 64-bit Ubuntu~24.04. It is installed
alongside the unchanged system Python. Both online and offline procedures create
\texttt{.venv\_ogs} and install OGS~6.5.9 and OGSTools~0.8.2.
\href{https://releases.ubuntu.com/}{Ubuntu~26.04} instead defaults to
Python~3.14 and requires \texttt{cp314} wheels. See the official
\href{https://ogstools.opengeosys.org/stable/user-guide/index.html}{OGSTools User Guide}
for further workflow information.

\section{Online installation on Ubuntu}

Check the operating system and architecture:

\begin{lstlisting}[style=terminal]
lsb_release -ds
uname -m
\end{lstlisting}

The architecture must be \texttt{x86\_64}. Ubuntu~24.04 provides Python~3.12
as its system Python, so Python~3.13 is installed in parallel from the
\href{https://launchpad.net/~deadsnakes/+archive/ubuntu/ppa}{deadsnakes PPA}.
Do not replace or remove Ubuntu's system Python. Install the repository tools,
add the PPA, and install Python~3.13 with its virtual-environment module:

\begin{lstlisting}[style=terminal]
sudo apt update
sudo apt install software-properties-common
sudo add-apt-repository ppa:deadsnakes/ppa
sudo apt update
sudo apt install python3.13 python3.13-venv
\end{lstlisting}

Create the course directory and environment, then install the pinned packages:

\begin{lstlisting}[style=terminal]
mkdir -p /path/to/OGS_Course
cd /path/to/OGS_Course
python3.13 -m venv .venv_ogs
source .venv_ogs/bin/activate
python -m pip install --upgrade pip setuptools wheel
python -m pip install \
  "ogs==6.5.9" "ogstools[all]==0.8.2" \
  "notebook" "jupyterlab"
python -c "from importlib.metadata import version; print('OGS:', version('ogs')); print('OGSTools:', version('ogstools'))"
\end{lstlisting}

The final command prints the installed OGS and OGSTools versions.

\section{Preparing the complete Ubuntu offline bundle}

Use an internet-connected Ubuntu computer with the \emph{same Ubuntu release
and CPU architecture} as the offline computer. Debian packages are specific to
the Ubuntu release; do not prepare them on Windows, macOS, or a different Ubuntu
release.

First download Python itself, pip, the virtual-environment module, and their
required Ubuntu packages:

\begin{lstlisting}[style=terminal]
mkdir -p /path/to/OGS_offline_bundle/python_debs
cd /path/to/OGS_offline_bundle
sudo apt update
sudo apt install software-properties-common apt-rdepends
sudo add-apt-repository ppa:deadsnakes/ppa
sudo apt update

apt-rdepends python3.13 python3.13-venv \
  | grep -v '^ ' \
  | grep -v '^<' \
  | sed 's/^debconf-2\.0$/debconf/' \
  | sort -u > python_packages.txt

cd python_debs
xargs -a /path/to/OGS_offline_bundle/python_packages.txt apt download
cd ..
\end{lstlisting}

The \texttt{sed} command replaces the virtual package name
\texttt{debconf-2.0} with its downloadable provider, \texttt{debconf}. The
\texttt{python\_debs} directory then contains Python~3.13 and its resolved
Ubuntu and deadsnakes dependencies. Any other package without a download
candidate indicates mismatched enabled repositories; resolve it before continuing.

Next download every Python wheel required by the course. Running this command
on matching Ubuntu avoids accidentally selecting wheels for another operating
system or architecture:

\begin{lstlisting}[style=terminal]
python3.13 -m venv .ogs_venv
source .ogs_venv/bin/activate
python -m pip install --upgrade pip
mkdir -p wheelhouse_ubuntu
python -m pip download --only-binary=:all: \
  -d wheelhouse_ubuntu \
  "ogs==6.5.9" "ogstools[all]==0.8.2" \
  "notebook" "jupyterlab"
deactivate
\end{lstlisting}

The download must complete without errors. Keep all downloaded wheel files;
removing an apparently unimportant dependency can make the offline installation
fail. Copy the course files into the bundle so that it has this layout:

\begin{lstlisting}[style=terminal]
OGS_offline_bundle/
|-- python_debs/
|   `-- Python and dependency .deb files
|-- wheelhouse_ubuntu/
|   `-- OGS, OGSTools, Jupyter, and dependency wheels
`-- OGS_Course/
    `-- course files
\end{lstlisting}

\section{Offline installation on Ubuntu}

Copy the complete bundle to the offline Ubuntu computer. Install the local
Python packages without contacting an online repository:

\begin{lstlisting}[style=terminal]
cd /path/to/OGS_offline_bundle/python_debs
sudo dpkg -i ./*.deb
sudo dpkg --configure -a
python3.13 --version
python3.13 -m venv --help
\end{lstlisting}

If \texttt{dpkg} reports a missing dependency, the bundle is incomplete. Return
to the matching online Ubuntu computer and download the named \texttt{.deb}; do
not replace it with a package from another Ubuntu release.

Create the environment and install only from the transferred wheelhouse:

\begin{lstlisting}[style=terminal]
cd /path/to/OGS_offline_bundle/OGS_Course
python3.13 -m venv .venv_ogs
source .venv_ogs/bin/activate
python -m pip install --no-index \
  --find-links /path/to/OGS_offline_bundle/wheelhouse_ubuntu \
  "ogs==6.5.9" "ogstools[all]==0.8.2" \
  "notebook" "jupyterlab"
python -c "from importlib.metadata import version; print('OGS:', version('ogs')); print('OGSTools:', version('ogstools'))"
\end{lstlisting}

The final command prints the installed OGS and OGSTools versions.

\chapter{OpenGeoSys Installation and Requirements on macOS}

The supplied macOS installer provides the same online and offline modes as the
Windows batch file. Because macOS does not execute Windows batch files, the Mac
equivalent is the executable shell script
\texttt{install\_ogstools\_macos.command}. It creates \texttt{.venv\_ogs}
beside the script and installs OGS~6.5.9, OGSTools~0.8.2, Jupyter Notebook, and
JupyterLab without modifying the system Python environment.

\section{Supported Mac configuration}

The binary wheel published for \texttt{ogs==6.5.9} and CPython~3.13 has the
platform tag \texttt{macosx\_26\_0\_arm64}. Therefore, this pinned course
installation requires:

\begin{itemize}
    \item an Apple Silicon Mac (M-series processor, reported as
          \texttt{arm64});
    \item macOS~26 or newer; and
    \item native ARM64 or universal2 Python~3.13.
\end{itemize}

There is no OGS~6.5.9 Intel macOS wheel. Intel Macs and Apple Silicon Macs
running an older macOS release cannot use this wheelhouse. For those computers,
use the course environment on Ubuntu or Windows, or use a separately validated
OGS version/build rather than silently changing the pinned course version.

Check the Mac before installation:

\begin{lstlisting}[style=terminal]
sw_vers -productVersion
uname -m
python3.13 --version
python3.13 -c "import platform; print(platform.machine())"
\end{lstlisting}

Both architecture commands must report \texttt{arm64}.

\section{Online installation on macOS}

An internet connection is required for this workflow. Download Python~3.13 from
\url{https://www.python.org/downloads/macos/}. Choose the signed macOS 64-bit
universal2 installer, run the \texttt{.pkg} file, and reopen Terminal. Verify it
with:

\begin{lstlisting}[style=terminal]
python3.13 --version
python3.13 -m pip --version
\end{lstlisting}

The universal2 installer supplies a native ARM64 Python when it is launched from
a native Terminal on Apple Silicon.

\subsection{Run the online macOS installer}

Open Terminal, change to the course directory, make the script executable, and
run it in online mode:

\begin{lstlisting}[style=terminal]
cd /path/to/OGS_Course_2026
chmod +x install_ogstools_macos.command
./install_ogstools_macos.command online
./.venv_ogs/bin/python -c "from importlib.metadata import version; print('OGS:', version('ogs')); print('OGSTools:', version('ogstools'))"
\end{lstlisting}

A \texttt{.command} file is an executable macOS shell script. Run
\texttt{chmod +x} only once; thereafter, double-clicking it in Finder or running
it without an argument selects online mode. Use the \texttt{offline} argument
shown below for offline mode. The script locates its own folder, keeps the
environment and wheelhouse beside itself, and verifies OGS, OGSTools, and Jupyter.

Activate the environment in a later Terminal session with:

\begin{lstlisting}[style=terminal]
cd /path/to/OGS_Course_2026
source .venv_ogs/bin/activate
\end{lstlisting}

\section{Offline installation on macOS}

First prepare the complete offline bundle as described in
Chapter~\ref{chap:macos-wheelhouse}. It must contain both the Python installer
itself and every Python wheel. Copy the bundle to the offline Mac with this
layout:

\begin{lstlisting}[style=terminal]
OGS_Course_2026/
|-- python-3.13.x-macos11.pkg
|-- install_ogstools_macos.command
`-- wheelhouse_macos_arm64/
    |-- ogs-6.5.9-...-macosx_26_0_arm64.whl
    |-- ogstools-0.8.2-...whl
    `-- all dependency wheel files
\end{lstlisting}

On the offline Mac, first double-click
\path{python-3.13.x-macos11.pkg} and complete the installation. Reopen Terminal
and verify Python before running the course installer:

\begin{lstlisting}[style=terminal]
python3.13 --version
python3.13 -m pip --version
\end{lstlisting}

Then open Terminal in the transferred directory and run:

\begin{lstlisting}[style=terminal]
chmod +x install_ogstools_macos.command
./install_ogstools_macos.command offline
./.venv_ogs/bin/python -c "from importlib.metadata import version; print('OGS:', version('ogs')); print('OGSTools:', version('ogstools'))"
\end{lstlisting}

Offline mode uses \texttt{--no-index} and
\texttt{--find-links wheelhouse\_macos\_arm64}; pip therefore does not contact
an online package index. The installation stops if the wheelhouse is missing or
incomplete.

\section{Using Jupyter on macOS}

Activate the environment and start either interface:

\begin{lstlisting}[style=terminal]
source .venv_ogs/bin/activate
jupyter notebook
# or
jupyter lab
\end{lstlisting}

To launch Notebook without activating the environment, use:

\begin{lstlisting}[style=terminal]
./.venv_ogs/bin/jupyter notebook
\end{lstlisting}

In VS Code, run \emph{Python: Select Interpreter} and choose
\path{OGS_Course_2026/.venv_ogs/bin/python}. Select the same interpreter as the
kernel when opening a Jupyter notebook.

\chapter{Preparing the Complete macOS Offline Bundle on Ubuntu}
\label{chap:macos-wheelhouse}

The complete bundle has two independent parts: the signed Python~3.13 macOS
installer and the Python wheelhouse. An Ubuntu computer can download macOS
wheels without installing them. pip must be told the target operating system,
CPU architecture, Python implementation, Python version, and ABI; otherwise it
downloads Linux wheels for the Ubuntu host.

\section{Download the Python installer itself}

On any internet-connected computer, open
\url{https://www.python.org/downloads/macos/} and download the signed Python~3.13
64-bit universal2 installer. The filename has the form
\path{python-3.13.x-macos11.pkg}. Store this \texttt{.pkg} file beside
\texttt{install\_ogstools\_macos.command}. The wheelhouse does not contain
Python itself, so this installer is mandatory for an offline Mac.

\section{Confirm the target tag on the Mac}

When the target Mac has temporary internet access, install the small
\texttt{packaging} utility and inspect its supported platform tags:

\begin{lstlisting}[style=terminal]
python3.13 -m pip install packaging
python3.13 -c "from packaging.tags import platform_tags; print(next(platform_tags()))"
\end{lstlisting}

For the supported course configuration, the result must include a platform tag
compatible with \path{macosx_26_0_arm64}.

\section{Download the Apple Silicon wheels from Ubuntu}

On the connected Ubuntu computer, update pip in a preparation environment and
download the complete CPython~3.13 wheel set:

\begin{lstlisting}[style=terminal]
python3 -m venv .ogs_venv
source .ogs_venv/bin/activate
python3 -m pip install --upgrade pip

mkdir -p wheelhouse_macos_arm64
python3 -m pip download --only-binary=:all: \
  --platform macosx_26_0_arm64 \
  --python-version 3.13 \
  --implementation cp --abi cp313 \
  -d wheelhouse_macos_arm64 \
  "ogs==6.5.9" "ogstools[all]==0.8.2" \
  "notebook" "jupyterlab"
\end{lstlisting}

The options select wheels for Apple Silicon macOS~26+, CPython~3.13, and the
\texttt{cp313} ABI. Platform-independent \texttt{py3-none-any} wheels are also
accepted. Do not replace the macOS platform with \texttt{linux\_x86\_64} or
copy Linux wheels into this directory.

The command must finish successfully. The binary-only option rejects source
archives. Consequently, a missing compatible wheel causes an error instead of
producing a bundle that would attempt compilation on the offline Mac.

\section{Verify and transfer the bundle}

Check that no source distributions or Linux/Windows wheels were collected:

\begin{lstlisting}[style=terminal]
find wheelhouse_macos_arm64 -maxdepth 1 -type f ! -name "*.whl"
find wheelhouse_macos_arm64 -maxdepth 1 -type f \
  \( -name "*linux*" -o -name "*win*" \)
ls wheelhouse_macos_arm64/ogs-6.5.9-*-macosx_26_0_arm64.whl
\end{lstlisting}

The first two commands should print nothing, and the last command should show
the OGS~6.5.9 CPython~3.13 macOS ARM64 wheel. Transfer the wheelhouse, the
macOS installer script, the course files, and the Python~3.13 universal2
\texttt{.pkg} installer to the offline Mac.

\chapter{OpenGeoSys Installation and Requirements on Windows}

The supplied batch installer supports both online and offline installation on
64-bit Windows. It creates an isolated Python environment beside the installer
and installs the course-tested versions OGS~6.5.9 and OGSTools~0.8.2. The
installer does not modify the system Python environment.

\section{Online installation on Windows}

Download and install the 64-bit Python~3.13 Windows installer from
\url{https://www.python.org/downloads/windows/}. During installation, enable
\emph{Add Python to PATH}, install the Python launcher, and ensure that
\texttt{pip} is selected. Open PowerShell or Command Prompt and verify the
installation:

\begin{lstlisting}[style=terminal]
py -3.13 --version
py -3.13 -m pip --version
\end{lstlisting}

Keep \texttt{install\_ogstools\_windows.bat} in the installation directory. The
script creates a virtual environment named \texttt{.venv\_ogs} in that directory.

\subsection{Run the online batch installer}

An internet connection is required for this mode. In PowerShell, change to the
directory containing the batch file and run:

\begin{lstlisting}[style=terminal]
.\install_ogstools_windows.bat online
.\.venv_ogs\Scripts\python.exe -c "from importlib.metadata import version; print('OGS:', version('ogs')); print('OGSTools:', version('ogstools'))"
\end{lstlisting}

A \texttt{.bat} file is a Windows command script. Run it in PowerShell or
Command Prompt with the \texttt{online} or \texttt{offline} argument;
double-clicking it or omitting the argument selects online mode. The script
locates its own folder, so the environment and wheelhouse remain beside it
regardless of the current working directory.

The script performs the following operations automatically:

\begin{enumerate}
    \item checks that Python~3.13 is available through the Windows launcher,
    \item creates or reuses \texttt{.venv\_ogs},
    \item installs the pinned OGS and OGSTools packages together with
          Jupyter Notebook and JupyterLab,
    \item checks the installed package versions and OGSTools configuration, and
    \item verifies that \texttt{ogs.exe} and Jupyter Notebook can be executed.
\end{enumerate}

A successful run ends with \texttt{Installation succeeded}. To use the
environment in a later PowerShell session, activate it with:

\begin{lstlisting}[style=terminal]
.\.venv_ogs\Scripts\Activate.ps1
\end{lstlisting}

For Command Prompt, use:

\begin{lstlisting}[style=terminal]
.\.venv_ogs\Scripts\activate.bat
\end{lstlisting}

If PowerShell prevents activation, allow locally created scripts for the current
user and then reopen PowerShell:

\begin{lstlisting}[style=terminal]
Set-ExecutionPolicy -ExecutionPolicy RemoteSigned -Scope CurrentUser
\end{lstlisting}

\section{Using Jupyter Notebook on Windows}
\label{subsec:windows-jupyter}

The batch installer places Jupyter Notebook and JupyterLab in the course virtual
environment, \texttt{.venv\_ogs}, together with OGS and OGSTools. This ensures that notebooks use
the course-tested Python packages instead of packages from the system Python.

Open PowerShell in the course directory and activate the environment:

\begin{lstlisting}[style=terminal]
Set-Location "C:\path\to\OGS_Course_2026"
.\.venv_ogs\Scripts\Activate.ps1
\end{lstlisting}

Start the traditional Jupyter Notebook interface with:

\begin{lstlisting}[style=terminal]
jupyter notebook
\end{lstlisting}

Alternatively, start it without activating the environment:

\begin{lstlisting}[style=terminal]
.\.venv_ogs\Scripts\jupyter.exe notebook
\end{lstlisting}

Jupyter normally opens the Notebook dashboard in the default web browser. The
folder from which it is started becomes the initial file browser location. Keep
the PowerShell window open while Jupyter is running, and press
\texttt{Ctrl+C} in that window to stop the server.

To use the newer JupyterLab interface instead, run:

\begin{lstlisting}[style=terminal]
jupyter lab
\end{lstlisting}

To confirm that a notebook uses the correct environment, create a Python
notebook and execute:

\begin{lstlisting}[style=terminal]
import sys
import ogstools
print(sys.executable)
print(ogstools.__version__)
\end{lstlisting}

The reported interpreter path should contain
\path{OGS_Course_2026\.venv_ogs\Scripts\python.exe}. These launch commands
are identical after online and offline installation.


\section{Installing and Configuring Visual Studio Code}
\label{subsec:windows-vscode}

Visual Studio Code (VS Code) is an optional graphical editor for the course.
It is installed separately from the Python environment because it is a Windows
application, not a Python package. The Python and Jupyter extensions connect
VS Code to the existing \texttt{.venv\_ogs} environment.

\subsection{Online installation}

Download the 64-bit Windows \emph{User Installer} from
\url{https://code.visualstudio.com/Download} and run it. The User Installer is
recommended for most participants because it normally does not require
administrator privileges. During setup, enable \emph{Add to PATH} if that
option is presented. Close and reopen PowerShell after installation, then
verify the command-line launcher:

\begin{lstlisting}[style=terminal]
code --version
\end{lstlisting}

Install the Microsoft Python and Jupyter extensions from PowerShell:

\begin{lstlisting}[style=terminal]
code --install-extension ms-python.python
code --install-extension ms-toolsai.jupyter
\end{lstlisting}

Open the course directory in VS Code:

\begin{lstlisting}[style=terminal]
Set-Location "C:\path\to\OGS_Course_2026"
code .
\end{lstlisting}

Press \texttt{Ctrl+Shift+P}, run \emph{Python: Select Interpreter}, and choose
the interpreter at
\path{OGS_Course_2026\.venv_ogs\Scripts\python.exe}. When opening an
\texttt{.ipynb} file, use the kernel selector in the upper-right corner of the
notebook editor and select the same \texttt{.venv\_ogs} environment. This makes
the terminal, Python files, and notebooks use the OGS course packages.

\subsection{Offline installation}

On an internet-connected Windows computer, download the VS Code User Installer
(or the ZIP archive when a portable installation is preferred). Also download
the \texttt{.vsix} packages for the Microsoft \emph{Python} and \emph{Jupyter}
extensions from their Visual Studio Marketplace pages. Copy the installer or
ZIP archive and all downloaded \texttt{.vsix} files to the offline computer.

Install VS Code, then install each extension from the VS Code Extensions view
using \emph{Views and More Actions} (the \texttt{...} menu) followed by
\emph{Install from VSIX}. If the \texttt{code} launcher is available, the same
operation can be performed in PowerShell:

\begin{lstlisting}[style=terminal]
code --install-extension .\vscode_extensions\ms-python.python-<version>.vsix
code --install-extension .\vscode_extensions\ms-toolsai.jupyter-<version>.vsix
\end{lstlisting}

If the Marketplace lists additional extension dependencies, download and
transfer their \texttt{.vsix} files as well. Finally, open the course folder and
select \path{.venv_ogs\Scripts\python.exe} as described above. The existing
offline Python installation and \texttt{wheelhouse\_windows} are unchanged;
VS Code only provides the editor and notebook interface.


\chapter{Preparing the Complete Windows Offline Bundle}
\label{chap:windows-offline-bundle}

Prepare the bundle on a computer with internet access. It must contain both the
Python installer and every required Python wheel.

First download the 64-bit Python~3.13 Windows installer from
\url{https://www.python.org/downloads/windows/}. Keep the downloaded
\path{python-3.13.x-amd64.exe} file; do not download the ARM64 or embeddable
package. The following commands can then be run on an online Ubuntu computer to
download binary wheels for 64-bit CPython~3.13 without installing them there:

\begin{lstlisting}[style=terminal]
mkdir -p wheelhouse_windows
python3 -m pip download --only-binary=:all: \
  --platform win_amd64 --python-version 3.13 \
  --implementation cp --abi cp313 \
  -d wheelhouse_windows \
  "ogs==6.5.9" "ogstools[all]==0.8.2" \
  "notebook" "jupyterlab"
\end{lstlisting}

The options have the following meanings:

\begin{itemize}
    \item \texttt{--only-binary=:all:} accepts wheels only and rejects source
          distributions that would require compilation;
    \item \texttt{--platform win\_amd64} selects 64-bit Windows packages;
    \item \texttt{--python-version 3.13} selects Python~3.13;
    \item \texttt{--implementation cp} selects CPython;
    \item \texttt{--abi cp313} selects the CPython~3.13 ABI; and
    \item \texttt{-d wheelhouse\_windows} stores the downloaded packages in the
          directory expected by the batch installer.
\end{itemize}

The download must finish successfully. If pip reports that no compatible binary
distribution exists, do not proceed with an incomplete wheelhouse. Prepare the
bundle with a recent pip version and keep every downloaded \texttt{.whl} file.

Copy the installer and the complete wheel directory to the offline Windows
computer. The required layout is:

\begin{lstlisting}[style=terminal]
OGS_Course_2026\
|-- python-3.13.x-amd64.exe
|-- install_ogstools_windows.bat
`-- wheelhouse_windows\
    |-- ogs-6.5.9-...whl
    |-- ogstools-0.8.2-...whl
    `-- all dependency wheel files
\end{lstlisting}

\chapter{Offline Installation on Windows}

Copy the complete bundle to the offline Windows computer. First double-click
\path{python-3.13.x-amd64.exe}. Enable \emph{Add Python to PATH}, install the
Python launcher, and keep \texttt{pip} selected. No internet connection is
required by the Python installer. Verify it in PowerShell:

\begin{lstlisting}[style=terminal]
py -3.13 --version
py -3.13 -m pip --version
\end{lstlisting}

Then open PowerShell or Command Prompt in the folder containing the batch file
and run:

\begin{lstlisting}[style=terminal]
.\install_ogstools_windows.bat offline
.\.venv_ogs\Scripts\python.exe -c "from importlib.metadata import version; print('OGS:', version('ogs')); print('OGSTools:', version('ogstools'))"
\end{lstlisting}

Offline mode uses:

\begin{lstlisting}[style=terminal]
pip install --no-index --find-links wheelhouse_windows ...
\end{lstlisting}

Consequently, pip does not contact an online package index. Installation fails
if the wheel directory is missing or if any required dependency is absent. The
same automatic package and executable checks used by online mode run after the
offline installation.

\section{Troubleshooting}

\begin{itemize}
    \item If \texttt{py -3.13} is not found, reinstall 64-bit Python~3.13 with
          the Python launcher enabled.
    \item If an earlier failed run left an unusable \texttt{.venv\_ogs}, remove only
          that \texttt{.venv\_ogs} directory and run the installer again.
    \item Do not mix wheels for Python~3.13, 32-bit Windows, or other operating
          systems with the CPython~3.13 64-bit Windows wheelhouse.
    \item Review the output immediately preceding the failure message. It normally
          identifies the missing package or failed check.
\end{itemize}

\end{document}

